\sscp{Glosses associated with *mantəR}{14-04-01}
In \srf{02-04-04} we noted that \citet{Blust1982,Blust2012b}
had translated \it{buidelrat} from Dutch sources as ‘bandicoot’,
parsing it literally as ‘pouch rat’.
However several modern Dutch dictionaries give English ‘possum,
opossum’ as the meaning associated with \it{buidelrat}.\footnote{
		For example, in the Dutch version of Wikipedia \it{buidelrat} redirects to
		\it{opossum} (Didelphidae) -- the family of marsupials found
		in the Americas. (‘Opossums’, Wikipedia. nl.wikipedia.org/wiki/Opossums.
		Accessed 22 June 2023)}
Possums are arboreal marsupials so close
in size, shape and behaviour to cuscus that they are often confused with
the cuscus where the two co-exist,
such as in northern Australia.
They both belong to the suborder Phalangeriformes which includes
“animals commonly known as possums,
gliders, and cuscus” (`Phalangeriformes', Wikipedia. en.wikipedia.org/wiki/Phalangeriformes.
Accessed 22 June 2023).

We noted that \citet[304--305]{vanEkris1864} provides independent
evidence via Malay that Dutch \it{buidelrat} was explicitly associated
with ‘cuscus’ during the era when many Dutch sources recorded
languages for which we have no other data.
For the Kamarian language of southwest Seram,
van Ekris gives \it{maker} ‘buidelrat’,
\it{make mauriat} ‘gevlekte buidelrat [Ml.
kusu potar]’ (Ml. = Malay), \it{make moon} ‘groote buidelrat’.
\it{Kusu} is the Ambon Malay term for ‘cuscus’.
And Dutch \it{gevlekte} = ‘spotted’,
so \it{make mauriat} ‘gevlekte buidelrat [Ml.
kusu potar]’ is explicitly identified in Malay as a type of \it{kusu}
and refers to the ‘spotted cuscus (\it{Spilocuscus maculatus})’.
This is clear evidence that the term \it{buidelrat} in 19th century
Dutch sources can refer to ‘cuscus’.

Similarly, \citet{Stresemann1927} used German terms associated with northern
hemisphere animals to describe animals of Seram.
His Proto-Ambon **maḏer is glossed interchangeably as both ‘Baummarder’ (e.g.
p.101) or ‘Beutelmarder’ (e.g. p. 129).
\it{Baummarder} = ‘marten, \it{Martes spp.}’ which is a nocturnal
tree-climbing animal, while \it{Beutelmarder} (literally ‘pouch marten’) is equated
by some modern sources with ‘quoll, \it{Dasyurus}’.
(e.g. ‘Beutelmarder’, Wikipedia. de.wikipedia.org/wiki/ Beutelmarder.
Accessed 22 June 2023).
Proto-Ambon **maḏer is the etymon for Paulohi \it{makel-e} (p. 41)
which is glossed in \citet[130]{Stresemann1918} as
``cuscus \it{Phalanger orientalis} and \it{Ph. maculatus}''.
Likewise, **maḏer is the etymon for
Nuaulu \it{maran-e} (p. 101) and Kamarian \it{maker} (p.
132), both of which are independently verified as ‘cuscus’ by
\citet[196f]{Ellen1972,Ellen1993} and \citet[304f]{vanEkris1864} respectively.
Furthermore, when the gloss of a reflex is different to his reconstruction
\citet{Stresemann1927} specifies the different gloss (such as on page 136
for Kobi: “Mond SS. \it{umala} > \it{umel-am} „Sonne”.” indicating semantic shift
from Proto-Sub-Seram **umala ‘moon’ to Kobi ‘sun’).
Therefore, all reflexes of Stresemann’s Proto-Ambon **maḏer can
be taken to mean cuscus.\footnote{
		We have identified only one case in which Stresemann gives a different
		gloss to a reflex of **maḏer. This is Apiolo \it{marel-e} which is
		specified as ``männl. Beutelmarder' = `male pouch marten' \citep[180]{Stresemann1927}.}

We also agree with \citet{Schapper2011} that \it{buidelrat} interpreted
as ‘bandicoot’ outside of or west of the region in which bandicoots
are known to inhabit is an unlikely interpretation.

There are additional challenges with the broad English gloss
‘marsupial’, such as that found in \citet{Collins1983} for many cognates of **mans(aə)r.
English ‘marsupial’ is a generic term that could refer to the
Phalanger cuscus, a tree kangaroo, a sugar glider, or a bandicoot.
But by extrapolation we have clear evidence that Collins uses
the gloss ‘marsupial’ to refer to the cuscus.
\citet[42]{Collins1981} gives the gloss ‘marsupial’ for the Buru word \it{tonal} ‘cuscus’.
The Buru word \it{tonal} refers to the cuscus
and not to any other animal (\cite[19]{Grimes1991c}, \cite[942]{GrimesGrimes2020}; see \frf{14-04}).
Furthermore tree kangaroos and bandicoots are not found as far west as Buru.
In the same article \citet[36]{Collins1981} glosses Ambelau \it{mate}
as ‘marsupial’, undoubtedly referring to the cuscus abundant on that island inhabited
by Muslims who do not hunt or eat it,
and using the same gloss he used for Buru \it{tonal} in the same publication.
\citet[241]{Blust1982} incorrectly reglossed Ambelau \it{mate}
as ‘bandicoot, marsupial rat’, apparently unaware that Ambelau is far from the
region where bandicoots are known to be found.
\citet{BlustTrussel2020} now accept that Ambelau \it{mate} should
be glossed ‘cuscus’.\footnote{
		\citet{BlustTrussel2020} incorrectly
		flag that the Ambelau term involves metathesis.}

\begin{figure}[h]\centering
	\includegraphics[width=0.6\columnwidth]{14-04.jpg}
	\fcp{Animal referred to as \it{tonal} by Buru speakers}{14-04}
\end{figure}

\citet[241]{Blust1982} also incorrectly glossed Asilulu \it{marel}
as ‘bandicoot, marsupial rat’, again unaware that Ambon is not in the region
where bandicoots are found,
with their known habitat not extending west of central Seram.
Regardless, \citet[70]{Collins2003a} provides the authoritative
gloss for Asilulu \it{marel} as “phalanger,
including [\it{Phalanger orientalis orientalis Pallas}], [\it{Ph.
maculatus chrysorrhous Temminek}] and [\it{Ph.
celebensis Gray}]”.\footnote{
		The Asilulu form with schwa \it{marəl}
		listed in \citet{BlustTrussel2020} is incorrect.
		Asilulu has no schwa
		and there is no metathesis as indicated by Blust
		and Trussel (Asilulu *nt/*nd/*ns > \it{r}; PMP *R > \it{l},
		as in *timuR > Asilulu \it{timul} ‘east’).}

Several of the words in other languages that \citet[55]{Collins1983}
etc.) glossed broadly as ‘marsupial’ are identified by
\citet{Ellen1972}[196f], \citeyear{Ellen1993} specifically as the Phalanger cuscus from data
that Ellen himself collected.

There is a practical problem working backwards as well.
What term was used by Collins
and others to elicit the words glossed with
the generic ‘marsupial’? For example,
the master Holle list \citep[50]{Stokhof1980a} equates item {\#}819
Dutch: \it{buideldier} `marsupial' (generic) = English: \it{marsupial} (generic)
= Indonesian: \it{kuskus} ‘cuscus’ (specific).
There is no generic term in Indonesian for ‘marsupial’ that is
known by rural people in eastern Indonesia, even educated ones.
The Indonesian term \it{kuskus} ‘cuscus’ is known in many regions.
In parts of Maluku,
Ambon Malay \it{kusu} ‘cuscus’ is better known.
Some also associate the Indonesian word \it{tupai} (glossed in
dictionaries as ‘squirrel, tree shrew’ -- both animals that do not exist in eastern Indonesia)
with their local terms for the marsupial cuscus; e.g.
\citet[77]{LaidigLaidig1991} gloss Larike \it{máridu} as: “tupai
[Melayu Ambon: kusu]; \it{k.o.
marsupial, cuscus}”,\footnote{
		Similarly, John Christensen (pers. comm.
		18 October 2020) and David Mead (pers. comm.
		October 2020), report the gloss ‘tupai’ from their consultants
		for words meaning ‘cuscus’.}
equating one arboreal mammal with another.

The preceding paragraphs all provide background for the basis
for correcting several glosses of English ‘marsupial’
and Dutch ‘buidelrat’ that have been incorrectly interpreted
as ‘bandicoot’, rather than ‘cuscus (possum, Phalanger)’.
This includes \citeauthor{Jonker1932}’s (\citeyear{Jonker1932}) disputed \it{buidelrat} for Leti-Moa \it{mada
{\tl} made}, which \citet[264]{Blust2012b} asserted must be ‘bandicoot’,
even though \citet[263]{Schapper2011} pointed out that \citet{vanEngelenhoven2004}
had listed Leti \it{mada} as ‘cuscus’.
And as we noted,
\it{buidelrat} is glossed as ‘possum’ in several Dutch dictionaries,
and was explicitly identified as ‘cuscus’ by \citet{vanEkris1864}.

